%% ============================================================
%%  Explicabilité en IA : de LIME à SoftCAM
%%  Panorama des méthodes post-hoc et intrinsèques
%%  Application à la prédiction de charge FaaS
%% ============================================================
\documentclass[10pt]{beamer}

\usepackage[utf8]{inputenc}
\usepackage[T1]{fontenc}
\usepackage[french]{babel}
\usepackage{lmodern}
\usepackage{graphicx}
\usepackage{booktabs}
\usepackage{array}
\usepackage{xcolor}
\usepackage{amsmath}
\usepackage{tikz}
\usetikzlibrary{arrows.meta, positioning, shapes.geometric, decorations.pathreplacing}

\usetheme{Madrid}
\usecolortheme{seahorse}

\usepackage[most]{tcolorbox}
\tcbset{colback=blue!5, colframe=blue!50!black, arc=3mm}
\newtcolorbox{dilemmabox}{
	colback=red!5,
	colframe=red!70!black,
	arc=3mm,
	fonttitle=\bfseries,
	title=Dilemme,
	halign title=center,
	halign=center
}

\newtcolorbox{constatbox}{
	colback=green!5,
	colframe=green!50!black,
	arc=3mm,
	fonttitle=\bfseries,
	title=Constat,
	halign title=center,
	halign=center
}

%% --- Couleurs personnalisées ---
\definecolor{shap}{RGB}{52,120,200}
\definecolor{cam}{RGB}{200,80,50}
\definecolor{softcam}{RGB}{40,160,80}
\definecolor{fayam}{RGB}{120,60,180}
\definecolor{verrou}{RGB}{180,30,30}

%% --- Commandes utilitaires ---
\newcommand{\source}[1]{\vspace{6pt}\hfill{\textcolor{gray}{\textbf{\tiny #1}}}}
\newcommand{\emphbox}[1]{%
  \begin{block}{}
    \centering\textbf{#1}
  \end{block}}

%% --- Métadonnées ---
\title{Méthodes d'Explicabilité en IA}
\subtitle{De LIME à SoftCAM --- Panorama et application à la prédiction de charge FaaS}
\author{FOSSO Cabrel}
\institute{Université de Dschang}
\date{\today}

%% --- Slide de transition automatique avant chaque section ---
\AtBeginSection[]{%
  \begin{frame}[plain, noframenumbering]
    \vfill
    \centering
    \begin{beamercolorbox}[sep=14pt, center, shadow=true, rounded=true]{title}
      \usebeamerfont{title}\insertsectionhead\par
    \end{beamercolorbox}
    \vspace{0.2cm}
    \tableofcontents[currentsection, hideallsubsections,
                     sectionstyle=show/shaded]
    \vfill
  \end{frame}%
}

%% ============================================================
\begin{document}
%% ============================================================

%% ------------------------------------------------------------------
%%  SECTION 0 — TITRE + SOMMAIRE
%% ------------------------------------------------------------------
\frame{\titlepage}

\begin{frame}{Plan de la présentation}
  \tableofcontents[hideallsubsections]
\end{frame}


%% ------------------------------------------------------------------
%%  SECTION 1 — L'EXPLICABILITÉ : POURQUOI MAINTENANT ?
%% ------------------------------------------------------------------
\section{L'explicabilité : pourquoi maintenant ?}

%% 1.1
\begin{frame}{Le paradoxe performances / confiance}
  \begin{columns}
    \begin{column}{0.55\textwidth}
      \begin{itemize}
        \item Les modèles ML modernes atteignent des performances remarquables :\\
              ResNet, GPT, Transformers \ldots
        \item Mais leur complexité les rend \alert{opaques} :
              \begin{itemize}
                \item Millions de paramètres
                \item Décisions non traçables
                \item Comportements inattendus hors distribution
              \end{itemize}
        \item Résultat : \textbf{les utilisateurs ne font pas confiance} aux prédictions
              qu'ils ne comprennent pas
      \end{itemize}
    \end{column}
    \begin{column}{0.42\textwidth}
    	\begin{dilemmabox}
    		Performance élevée\\[4pt]
    		\textbf{vs}\\[4pt]
    		Confiance \& déployabilité
    	\end{dilemmabox}
    	
    	\vspace{6pt}
    	
    	\begin{constatbox}
    		Un modèle précis mais inexplicable\\
    		\textbf{ne peut pas être déployé}\\
    		dans un secteur critique
    	\end{constatbox}
    \end{column}
  \end{columns}
\end{frame}

%% 1.2
\begin{frame}{Enjeux concrets : réglementation et secteurs critiques}
	\begin{tcolorbox}[title=Réglementation]
		\begin{itemize}
			\item \textbf{EU AI Act (2024)} : exige la traçabilité des décisions automatisées
			dans les systèmes à risque élevé
			\item \textbf{RGPD} : droit à l'explication pour les décisions automatisées
		\end{itemize}
	\end{tcolorbox}
	
	\vspace{6pt}
	
	\begin{tabular}{p{0.3\textwidth}p{0.3\textwidth}p{0.3\textwidth}}
		\hline
		\textbf{Santé} & \textbf{Finance} & \textbf{Cloud / FaaS} \\
		\hline
		Diagnostic, imagerie médicale : le médecin \textbf{doit comprendre} la prédiction &
		Décision de crédit, détection de fraude : obligation légale d'explication &
		Allocation de ressources : les opérateurs ont besoin de \textbf{comprendre les pics de charge} \\
		\hline
	\end{tabular}
	
	\vspace{8pt}
	\source{EU AI Act, 2024 ; RGPD Art. 22}
\end{frame}


%% 1.3
\begin{frame}{Qu'est-ce qu'une bonne explication ?}
	\begin{block}{3 dimensions clés}
		\begin{itemize}
			\item \textbf{Portée}
			\begin{itemize}
				\item \alert{Locale} : expliquer \textit{une} prédiction
				\item \alert{Globale} : expliquer le comportement général du modèle
			\end{itemize}
			\item \textbf{Fidélité}\\
			L'explication reflète-t-elle vraiment le raisonnement du modèle ?
			\item \textbf{Interprétabilité}\\
			L'explication est-elle compréhensible par un humain non expert ?
		\end{itemize}
	\end{block}
	
	\vspace{8pt}
	
	\begin{exampleblock}{3 propriétés souhaitables (SHAP)}
		\begin{itemize}
			\item \textbf{Précision locale} : l'explication prédit correctement localement
			\item \textbf{Missingness} : une feature absente a contribution nulle
			\item \textbf{Cohérence} : si une feature contribue plus dans un modèle,
			son score ne doit pas diminuer
		\end{itemize}
	\end{exampleblock}
\end{frame}


%% ------------------------------------------------------------------
%%  SECTION 2 — CONTEXTE : LES TRAVAUX DE FAYAM
%% ------------------------------------------------------------------
\section{Contexte : les travaux de FAYAM}

%% 2.1
\begin{frame}{Prédiction de charge FaaS : le problème métier}
  \begin{columns}
    \begin{column}{0.55\textwidth}
      \begin{itemize}
        \item \textbf{FaaS} (Function-as-a-Service) : paradigme cloud où les fonctions
              s'exécutent à la demande
        \item \textbf{Défi} : les invocations sont imprévisibles ---
              pics de charge soudains, saisonnalités complexes
        \item \textbf{Enjeu} : pré-allouer les ressources pour éviter la latence
              ou le gaspillage
        \item \textbf{Solution} : prédire le nombre d'invocations futures
              à partir de l'historique
      \end{itemize}
    \end{column}
    \begin{column}{0.42\textwidth}
      \begin{tcolorbox}[
      	title=Dataset FAYAM,
      	colback=blue!5,
      	colframe=blue!50!black,
      	arc=3mm,
      	fonttitle=\bfseries
      	]
      	\begin{itemize}
      		\item 18 fonctions FaaS instrumentées
      		\item Séries temporelles multivariées
      		\item 33 clusters de profils de charge
      		(DBSCAN --- k-distance)
      		\item Horizon de prédiction : \textit{predictionLength}
      	\end{itemize}
      \end{tcolorbox}
    \end{column}
  \end{columns}
\end{frame}

%% 2.2
\begin{frame}{Les deux architectures proposées par FAYAM}
	\textbf{Modèle 1 --- CNN-LSTM}
	\rule{\linewidth}{0.5pt}
	\vspace{2pt}
	
	\begin{itemize}
		\item Extraction de features locales (CNN)
		\item Mémorisation séquentielle (LSTM)
		\item Approche hybride
	\end{itemize}
	
	{\small \textit{(Mentionné pour exhaustivité --- hors de notre périmètre)}}
	
	\vspace{6pt}
	
	\textbf{Modèle 2 --- Transformer} \textcolor{blue!70!black}{\textbf{[focus]}}
	\rule{\linewidth}{0.5pt}
	\vspace{2pt}
	
	\begin{itemize}
		\item \texttt{TimeSeriesTransformer} HuggingFace
		\item Encodeur--Décodeur
		\item Mécanisme d'attention multi-têtes
		\item \texttt{output\_attentions=True}
	\end{itemize}
	
	{\small \textcolor{blue!70!black}{C'est ce modèle que nous allons chercher à expliquer}}
\end{frame}

%% 2.3
\begin{frame}{Résultats du Transformer sur les datasets de test}
	\begin{tcolorbox}[
		colback=blue!3,
		colframe=blue!40!black,
		arc=3mm,
		boxrule=0.8pt,
		fonttitle=\bfseries,
		title={Métriques de performance (phase test --- FAYAM, 2024)}
		]
		\begin{tabular}{lcccc}
			\toprule
			\textbf{Modèle} & \textbf{sMAPE} & \textbf{RMSE} & \textbf{R²} & \textbf{Spearman} \\
			\midrule
			CNN-LSTM        & \textit{(valeur)} & \textit{(valeur)} & \textit{(valeur)} & \textit{(valeur)} \\
			\textbf{Transformer} & \textit{(valeur)} & \textit{(valeur)} & \textit{(valeur)} & \textit{(valeur)} \\
			\bottomrule
		\end{tabular}
		
		\vspace{6pt}
		{\small \textit{Note : à compléter après reproduction des expériences FAYAM (Phase 1 du projet)}}
	\end{tcolorbox}
	
	\vspace{10pt}
	
	\begin{tcolorbox}[
		colback=gray!3,
		colframe=gray!50!black,
		arc=3mm,
		boxrule=0.8pt,
		fonttitle=\bfseries,
		title={Interprétation}
		]
		\begin{itemize}
			\item Les performances sont \textbf{compétitives} par rapport à l'état de l'art FaaS
			\item Le Transformer surpasse le CNN-LSTM sur la majorité des datasets
		\end{itemize}
	\end{tcolorbox}
\end{frame}

%% 2.4
\begin{frame}{Le verrou : de bonnes performances, mais aucune explication}
	\begin{center}
		\begin{tikzpicture}[
			node distance=1.2cm,
			box/.style={draw, rounded corners, text width=3.8cm, align=center, minimum height=0.9cm},
			redbox/.style={draw=verrou, fill=verrou!10, rounded corners, text width=5.5cm, align=center, minimum height=1.1cm},
			]
			\node[box, fill=fayam!15] (model) {\textbf{Transformer FAYAM}\\(bien entraîné)};
			\node[box, fill=green!15, right=of model] (pred) {\textbf{Prédiction}\\sMAPE faible};
			\node[redbox, below=of pred] (verrou) {\alert{\textbf{Aucune explication associée}}\\Pourquoi cette prédiction ?\\Quelles features ont compté ?};
			\draw[-{Stealth}] (model) -- (pred);
			\draw[-{Stealth}, dashed, verrou] (pred) -- (verrou);
		\end{tikzpicture}
	\end{center}
	
	\vspace{4pt}
	
	\begin{tcolorbox}[
		colback=red!3,
		colframe=red!60!black,
		arc=3mm,
		boxrule=0.8pt,
		fonttitle=\bfseries,
		title={Conséquence directe}
		]
		Un opérateur FaaS ne peut pas \textbf{valider}, \textbf{auditer} ni \textbf{faire confiance}
		à un modèle qui ne justifie pas ses prédictions.
		
		\vspace{6pt}
		\centering
		\textbf{La déployabilité dans un secteur critique : impossible}
	\end{tcolorbox}
\end{frame}


%% ------------------------------------------------------------------
%%  SECTION 3 — CARTE DES MÉTHODES D'EXPLICABILITÉ
%% ------------------------------------------------------------------
\section{Taxonomie des méthodes d'explicabilité}

%% 3.1
\begin{frame}{Axes de classification}
	\begin{columns}
		\begin{column}{0.48\textwidth}
			\begin{tcolorbox}[
				colback=blue!3,
				colframe=blue!40!black,
				arc=2mm,
				boxrule=0.5pt,
				fonttitle=\bfseries,
				title={Axe 1 --- Quand ?}
				]
				\begin{itemize}
					\item \textbf{Post-hoc} : \textit{(LIME, SHAP, TimeSHAP)}
					\item \textbf{Ante-hoc} : \textit{(arbres décision, régression log.)}
					\item \textbf{Pendant l'entraînement} : \textit{(SoftCAM $\leftarrow$ focus)}
				\end{itemize}
			\end{tcolorbox}
		\end{column}
		\begin{column}{0.48\textwidth}
			\begin{tcolorbox}[
				colback=green!3,
				colframe=green!40!black,
				arc=2mm,
				boxrule=0.5pt,
				fonttitle=\bfseries,
				title={Axe 2 --- Qui ?}
				]
				\begin{itemize}
					\item \textbf{Agnostique} : \textit{(LIME, SHAP)}
					\item \textbf{Spécifique} : \textit{(CAM, GradCAM, SoftCAM)}
				\end{itemize}
			\end{tcolorbox}
		\end{column}
	\end{columns}
	
	\vspace{10pt}
	
	\begin{columns}
		\begin{column}{0.48\textwidth}
			\begin{tcolorbox}[
				colback=orange!3,
				colframe=orange!50!black,
				arc=2mm,
				boxrule=0.5pt,
				fonttitle=\bfseries,
				title={Axe 3 --- Quelle portée ?}
				]
				\begin{itemize}
					\item \textbf{Locale}
					\item \textbf{Globale}
					\item \textbf{Semi-locale} : \textit{(TsSHAP)}
				\end{itemize}
			\end{tcolorbox}
		\end{column}
		\begin{column}{0.48\textwidth}
			\begin{tcolorbox}[
				colback=purple!3,
				colframe=purple!50!black,
				arc=2mm,
				boxrule=0.5pt,
				fonttitle=\bfseries,
				title={Axe 4 --- Quelle tâche ?}
				]
				\begin{itemize}
					\item \textbf{Classification}
					\item \textbf{Régression}
					\item \textbf{Prévision} : \textit{$\leftarrow$ notre cas FaaS}
				\end{itemize}
			\end{tcolorbox}
		\end{column}
	\end{columns}
\end{frame}

%% 3.2
\begin{frame}{Positionnement de toutes les méthodes}
	\begin{tcolorbox}[
		colback=blue!3,
		colframe=blue!40!black,
		arc=3mm,
		boxrule=0.8pt,
		enhanced,
		fonttitle=\bfseries,
		title=Comparaison des méthodes d'explicabilité
		]
		\begin{center}
			\scriptsize
			\begin{tabular}{lccccc}
				\toprule
				\textbf{Méthode} & \textbf{Quand ?} & \textbf{Agnostic ?} & \textbf{Prévision ?} & \textbf{Transformer ?} & \textbf{Famille} \\
				\midrule
				LIME             & Post-hoc & \checkmark & \texttimes & \checkmark & Perturbation \\
				KernelSHAP       & Post-hoc & \checkmark & \texttimes & \checkmark & SHAP \\
				TimeSHAP         & Post-hoc & \checkmark & \texttimes & partiel    & SHAP \\
				WindowSHAP       & Post-hoc & \checkmark & \texttimes & partiel    & SHAP \\
				TsSHAP           & Post-hoc & \checkmark & \checkmark & \checkmark & SHAP \\
				SHAPformer       & Post-hoc & partiel    & \checkmark & \checkmark & SHAP \\
				\midrule
				CAM              & Post-hoc & \texttimes & \texttimes & \texttimes & CAM \\
				GradCAM          & Post-hoc & \texttimes & \texttimes & \texttimes & CAM \\
				\midrule
				\textbf{SoftCAM} & \textbf{Entraîn.} & \texttimes & \textbf{\checkmark*} & \textbf{\checkmark*} & \textbf{CAM} \\
				\bottomrule
			\end{tabular}
		\end{center}
	\end{tcolorbox}

		\centering
		\footnotesize
		\textcolor{blue!70!black}{*} Extension proposée dans notre projet (H1) : transposition du principe SoftCAM au Transformer de FAYAM
	
	\begin{tcolorbox}[
		colback=green!3,
		colframe=green!50!black,
		arc=3mm,
		boxrule=0.8pt,
		fonttitle=\bfseries,
		title={Notre positionnement}
		]
		\centering
		Aucune méthode existante ne couvre parfaitement notre cas ---
		\textbf{c'est précisément notre contribution}
	\end{tcolorbox}
\end{frame}


%% ------------------------------------------------------------------
%%  SECTION 4 — FAMILLE SHAP
%% ------------------------------------------------------------------
\section{Famille SHAP : des perturbations locales aux valeurs exactes}

%% --- Sous-section A : Fondations ---

%% 4.1a — LIME
\begin{frame}{LIME --- Ribeiro, Singh \& Guestrin (KDD 2016)}
	\begin{minipage}{0.49\textwidth}
		\begin{tcolorbox}[
			colback=blue!3,
			colframe=blue!40!black,
			arc=2mm,
			boxrule=0.5pt,
			fonttitle=\bfseries,
			title={Problème résolu}
			]
			Comment expliquer la prédiction d'un modèle quelconque
			de façon \textbf{localement fidèle} et \textbf{interprétable} ?
		\end{tcolorbox}
		
		\begin{tcolorbox}[
			colback=green!3,
			colframe=green!40!black,
			arc=2mm,
			boxrule=0.5pt,
			fonttitle=\bfseries,
			title={Méthode}
			]
			Apprendre un modèle linéaire sparse $g$ autour de $x$ :
			\vspace{-0.2cm}
			\[
			\xi(x) = \arg\min_{g \in G}\, \mathcal{L}(f, g, \pi_x) + \Omega(g)
			\]
			\vspace{-0.5cm}
			\begin{itemize}
				\item $\pi_x(z) = \exp(-D(x,z)^2/\sigma^2)$ : noyau de proximité
				\item $\Omega(g)$ : complexité (nb de features non nulles)
			\end{itemize}
		\end{tcolorbox}
	\end{minipage}
	\hfill
	\begin{minipage}{0.48\textwidth}
		\begin{tcolorbox}[
			colback=orange!3,
			colframe=orange!50!black,
			arc=2mm,
			boxrule=0.5pt,
			fonttitle=\bfseries,
			title={Représentations interprétables}
			]
			\begin{itemize}
				\item Texte : présence ou absence de mots
				\item Image : super-pixels (patches)
				\item Tabulaire : présence ou absence de features
			\end{itemize}
		\end{tcolorbox}
		
		\begin{tcolorbox}[
			colback=red!3,
			colframe=red!60!black,
			arc=2mm,
			boxrule=0.5pt,
			fonttitle=\bfseries,
			title={Algorithme (K-Lasso)}
			]
			\begin{enumerate}
				\item Échantillonner $N$ instances autour de $x$
				\item Pondérer par $\pi_x$
				\item Lasso → garder $K$ features
				\item OLS → poids finaux
			\end{enumerate}
		\end{tcolorbox}
	\end{minipage}
	
	\vspace{8pt}
	\source{Ribeiro et al., KDD 2016 --- arXiv:1602.04938}
\end{frame}

%% 4.1b — LIME résultats + SP-LIME
\begin{frame}{LIME --- Résultats et SP-LIME}
	\begin{minipage}{0.5\textwidth}
		\begin{tcolorbox}[
			colback=blue!3,
			colframe=blue!40!black,
			arc=2mm,
			boxrule=0.5pt,
			fonttitle=\bfseries,
			title={Résultats (fidélité)}
			]
			Recall sur features importantes :
			\begin{itemize}
				\item \textbf{LIME : > 90\%} sur tous les classifieurs (books + DVDs)
				\item Greedy : 64\% (books) --- 33\% (DVDs)
				\item Random : $\sim$17\%
			\end{itemize}
			F1 trustworthiness LIME : \textbf{96,6 / 94,5 / 96,2 / 96,7}\\
			(LR / NN / RF / SVM sur Books)
		\end{tcolorbox}
		
		\vspace{-0.3cm}
		
		\begin{tcolorbox}[
			colback=green!3,
			colframe=green!40!black,
			arc=2mm,
			boxrule=0.5pt,
			fonttitle=\bfseries,
			title={SP-LIME --- confiance globale}
			]
			Sélectionner $B$ instances représentatives et \textbf{non-redondantes}
			par optimisation sous-modulaire greedy
		\end{tcolorbox}
	\end{minipage}
	\hfill
	\begin{minipage}{0.49\textwidth}
		\begin{tcolorbox}[
			colback=red!3,
			colframe=red!60!black,
			arc=2mm,
			boxrule=0.5pt,
			fonttitle=\bfseries,
			title={Limite clé}
			]
			\begin{itemize}
				\item \textbf{Instabilité} : deux runs = deux explications différentes
				\item \textbf{Fidélité locale $\neq$ fidélité globale}
				\item Pas de garanties théoriques sur cohérence et missingness
				\item \textbf{Pas adapté aux séries temporelles} : features traitées comme indépendantes
			\end{itemize}
		\end{tcolorbox}
		
		\begin{tcolorbox}[
			colback=orange!3,
			colframe=orange!50!black,
			arc=2mm,
			boxrule=0.5pt,
			fonttitle=\bfseries,
			title={Exemple célèbre}
			]
			Husky vs Wolf : LIME révèle que le classifieur utilise la \textbf{neige en arrière-plan}, pas l'animal
		\end{tcolorbox}
	\end{minipage}
	
	\vspace{8pt}
	\source{Ribeiro et al., KDD 2016}
\end{frame}

%% 4.1c — LIME limites et transition
\begin{frame}{LIME --- Ce que cela nous enseigne}
  \begin{columns}
    \begin{column}{0.56\textwidth}
      \begin{tcolorbox}[
        colback=green!3, colframe=green!40!black,
        arc=2mm, boxrule=0.5pt, fonttitle=\bfseries,
        title={Apports fondamentaux}]
        \begin{itemize}
          \item Premier cadre \textbf{model-agnostique} rigoureux pour l'XAI post-hoc
          \item Notion de \textbf{représentation interprétable}
          \item SP-LIME : pont vers la confiance globale
        \end{itemize}
      \end{tcolorbox}
    \end{column}
    \begin{column}{0.4\textwidth}
      \begin{tcolorbox}[
        colback=red!3, colframe=red!60!black,
        arc=2mm, boxrule=0.5pt, fonttitle=\bfseries,
        title={Ce qui manque}]
        \begin{itemize}
          \item Garanties théoriques solides
          \item Stabilité des explications
          \item Adaptation au domaine temporel
        \end{itemize}
      \end{tcolorbox}
    \end{column}
  \end{columns}

  \vspace{6pt}
  \begin{tcolorbox}[colback=blue!3, colframe=blue!40!black, arc=2mm, boxrule=0.5pt]
    \centering
    Un an plus tard (2017), Lundberg \& Lee unifieront LIME et d'autres méthodes
    sous un cadre théorique unique basé sur la \textbf{théorie des jeux coopératifs}\\[4pt]
    La naissance de \textbf{SHAP --- Shapley Additive Explanations}
  \end{tcolorbox}
\end{frame}

%% 4.2a — KernelSHAP
\begin{frame}{KernelSHAP --- Lundberg \& Lee (NeurIPS 2017)}
  \begin{columns}
    \begin{column}{0.5795\textwidth}
    \vspace{-0.5cm}
      \begin{tcolorbox}[
        colback=blue!3, colframe=blue!40!black,
        arc=2mm, boxrule=0.5pt, fonttitle=\bfseries,
        title={La valeur de Shapley}]
        Issue de la théorie des jeux coopératifs (Shapley, 1953).
        Contribution \textbf{juste et unique} de chaque feature $i$ :
       {\footnotesize
       	\[
          \phi_i = \sum_{S \subseteq F \setminus \{i\}} \frac{|S|!(|F|-|S|-1)!}{|F|!}
          \left[f(S \cup \{i\}) - f(S)\right]
        \]}
        $\Rightarrow$ contribution marginale moyennée sur toutes les coalitions
      \end{tcolorbox}
      \begin{tcolorbox}[
        colback=orange!3, colframe=orange!50!black,
        arc=2mm, boxrule=0.5pt, fonttitle=\bfseries,
        title={KernelSHAP}]
        Reformule le calcul des valeurs de Shapley comme une régression
        pondérée (LASSO) --- approximation de LIME avec le \textbf{bon noyau}
      \end{tcolorbox}
    \end{column}
    \begin{column}{0.42\textwidth}
    	\vspace{-0.7cm}
      \begin{tcolorbox}[
        colback=green!3, colframe=green!40!black,
        arc=2mm, boxrule=0.5pt, fonttitle=\bfseries,
        title={SHAP unifie LIME, DeepLIFT, SHAP linéaire\ldots}]
        Seule solution satisfaisant \textbf{simultanément} :
        \begin{enumerate}
          \item \textbf{Précision locale}
          \item \textbf{Missingness}
          \item \textbf{Cohérence}
        \end{enumerate}
      \end{tcolorbox}
      \begin{tcolorbox}[
        colback=red!3, colframe=red!60!black,
        arc=2mm, boxrule=0.5pt, fonttitle=\bfseries,
        title={Limite}]
        {\small
        Calcul exact : $2^{|F|}$ coalitions.\\
        KernelSHAP \textbf{approxime} par échantillonnage.
        Coût élevé pour longues séquences.}
      \end{tcolorbox}
    \end{column}
  \end{columns}
  \source{Lundberg \& Lee, NeurIPS 2017}
\end{frame}

%% 4.2b — KernelSHAP : le problème séquentiel
\begin{frame}{KernelSHAP face aux séquences temporelles}
  \begin{tcolorbox}[
    colback=red!3, colframe=red!60!black,
    arc=2mm, boxrule=0.5pt, fonttitle=\bfseries,
    title={Le problème central}]
    KernelSHAP attribue de l'importance aux features du \textbf{pas de temps courant uniquement}.\\
    Il ignore l'historique de la séquence et l'état caché du modèle récurrent.
  \end{tcolorbox}

  \begin{center}
    \begin{tikzpicture}[scale=0.75, transform shape]
      \node[draw, fill=gray!20, minimum width=1.2cm, minimum height=0.7cm] (e1) at (0,0) {$e_1$};
      \node[draw, fill=gray!20, minimum width=1.2cm, minimum height=0.7cm] (e2) at (1.5,0) {$e_2$};
      \node[draw, fill=gray!20, minimum width=1.2cm, minimum height=0.7cm] (en) at (3.5,0) {$e_{l-1}$};
      \node[draw, fill=shap!40, minimum width=1.2cm, minimum height=0.7cm] (el) at (5,0) {$e_l$};
      \node[above=0.4cm of el, align=center, text=shap] {\small SHAP\\classique};
      \draw[-{Stealth}] (e1) -- (e2);
      \draw[dashed] (e2) -- (en);
      \draw[-{Stealth}] (en) -- (el);
      \draw[-{Stealth}] (el) -- ++(1.2,0) node[right] {Prédiction};
      \draw[thick, red!60, decorate, decoration={brace, amplitude=5pt}]
        (-0.72,0.5) -- (4.2,0.5) node[midway, above=5pt, red!60] {\small Ignoré par KernelSHAP};
    \end{tikzpicture}
  \end{center}

  \begin{tcolorbox}[colback=blue!3, colframe=blue!40!black, arc=2mm, boxrule=0.5pt]
    \centering
    $\Rightarrow$ Il faut adapter SHAP au domaine séquentiel : \textbf{TimeSHAP}
  \end{tcolorbox}
\end{frame}

%% --- PAUSE : les 3 propriétés ---
\begin{frame}{\textbf{Pause} --- Pourquoi SHAP s'impose face à LIME}
  \begin{columns}
    \begin{column}{0.466\textwidth}
      \begin{tcolorbox}[
        colback=red!3, colframe=red!60!black,
        arc=2mm, boxrule=0.5pt, fonttitle=\bfseries,
        title={LIME}]
        \begin{itemize}
          \item[\checkmark] Précision locale
          \item[$\times$] Missingness (non garantie)
          \item[$\times$] Cohérence (non garantie)
          \item[$\times$] Instabilité due au sampling
          \item[$\times$] Pas de fondement axiomatique
        \end{itemize}
      \end{tcolorbox}
    \end{column}
    \begin{column}{0.47\textwidth}
      \begin{tcolorbox}[
        colback=green!3, colframe=green!40!black,
        arc=2mm, boxrule=0.5pt, fonttitle=\bfseries,
        title={SHAP (KernelSHAP)}]
        \begin{itemize}
          \item[\checkmark] Précision locale
          \item[\checkmark] Missingness
          \item[\checkmark] Cohérence
          \item[\checkmark] Solution \textbf{unique} par les axiomes de Shapley
          \item[\checkmark] TreeSHAP exact en temps polynomial
        \end{itemize}
      \end{tcolorbox}
    \end{column}
  \end{columns}

  \vspace{6pt}
  \begin{tcolorbox}[
    colback=blue!3, colframe=blue!40!black,
    arc=2mm, boxrule=0.5pt, fonttitle=\bfseries,
    title={Conclusion de cette pause}]
    \centering
    La famille SHAP est le \textbf{standard théorique} pour l'XAI post-hoc.\\
    Les prochaines méthodes sont toutes des \textbf{extensions SHAP}\\
    adaptées aux séries temporelles.
  \end{tcolorbox}
\end{frame}

%% --- Sous-section B : Extensions séries temporelles ---

%% 4.4a — TimeSHAP
\begin{frame}{TimeSHAP --- Bento et al. (KDD 2021)}
	\textbf{Problème résolu}
	\rule{\linewidth}{0.5pt}
	
	SHAP ignore l'historique des RNN. TimeSHAP étend KernelSHAP
	au domaine séquentiel avec des \textbf{perturbations bidimensionnelles}.
	
	\vspace{8pt}
	
	\textbf{Méthode}
	\rule{\linewidth}{0.5pt}
	
	Matrice $X \in \mathbb{R}^{d \times l}$ ($d$ features, $l$ événements)
	\begin{itemize}
		\item \textbf{Features} : $h_X^f(z) = D_z X + (I-D_z)B$
		\item \textbf{Événements} : $h_X^e(z) = XD_z + B(I-D_z)$
		\item \textbf{Cellules} : intersection features $\times$ événements
	\end{itemize}
	
	\vspace{8pt}
	
	\textbf{Élagage temporel}
	\rule{\linewidth}{0.5pt}
	
	$O(2^l) \to O(l \cdot 2^{l-i})$ \quad
	Médiane : $138{,}5 \to \mathbf{14{,}0}$ événements ($\eta=0{,}025$)
	
	\vspace{6pt}
	\source{Bento et al., KDD 2021 --- arXiv:2012.00073}
\end{frame}

\begin{frame}{TimeSHAP --- Résultats et limites}
	\textbf{Résultats (fraude bancaire)}
	\rule{\linewidth}{0.5pt}
	\vspace{2pt}
	
	\begin{itemize}
		\item GRU + feed-forward, $\sim$20M instances
		\item Recall 84,3\% @ 1\% FPR
		\item Événement $t=0$ : \textbf{41\%} du score
		\item Pattern : enrollment-login-transaction
		\item Biais détecté : âge $\leftrightarrow$ taux FP
	\end{itemize}
	
	\vspace{8pt}
	
	\textbf{Limite}
	\rule{\linewidth}{0.5pt}
	\vspace{2pt}
	
	Conçu pour la \textbf{classification}. \\
	\textcolor{red}{Pas adapté à la prévision (FAYAM).}
	
	\vspace{6pt}
	\source{Bento et al., KDD 2021 --- arXiv:2012.00073}
\end{frame}

%% 4.4b
\begin{frame}{TimeSHAP --- Élagage temporel et explications cellules}
  \begin{columns}
    \begin{column}{0.569\textwidth}
    \vspace{-1.5cm}
      \begin{tcolorbox}[
        colback=orange!3, colframe=orange!50!black,
        arc=2mm, boxrule=0.5pt, fonttitle=\bfseries,
        title={Algorithme d'élagage (Alg. 1)}]
        \begin{enumerate}
          \item Parcourir la séquence de la fin vers le début
          \item Si $|w_1| < \eta$ : grouper les événements passés
          \item Complexité finale : $O(l \cdot 2^{l-i})$
        \end{enumerate}
      \end{tcolorbox}
      \begin{tcolorbox}[
        colback=blue!3, colframe=blue!40!black,
        arc=2mm, boxrule=0.5pt, fonttitle=\bfseries,
        title={Pruning Table (Table 1)}]
        \hspace{-0.5cm}
        \begin{tabular}{lcc}
          \toprule
          & \textbf{Original} & $\eta=0{,}025$ \\
          \midrule
          Longueur médiane & 138,5 & \textbf{14,0} \\
          RSD & 1,71 & \textbf{0,98} \\
          \bottomrule
        \end{tabular}
      \end{tcolorbox}
    \end{column}
    \begin{column}{0.43\textwidth}
    \vspace{-0.7cm}
      \begin{tcolorbox}[
        colback=green!3, colframe=green!40!black,
        arc=2mm, boxrule=0.5pt, fonttitle=\bfseries,
        title={3 niveaux d'explication}]
        \begin{itemize}
          \item \textbf{Événements} : quels instants ?
          \item \textbf{Features} : quelles variables ?
          \item \textbf{Cellules} : quels (feature, instant) ?
        \end{itemize}
      \end{tcolorbox}
      \begin{tcolorbox}[
        colback=red!3, colframe=red!60!black,
        arc=2mm, boxrule=0.5pt, fonttitle=\bfseries,
        title={Transposabilité à FAYAM}]
        S'applique en boîte noire au Transformer FAYAM, mais explique
        un score de probabilité --- pas une prévision de charge.
        Adaptation requise.
      \end{tcolorbox}
    \end{column}
  \end{columns}
  \source{Bento et al., KDD 2021}
\end{frame}

%% 4.5 — WindowSHAP
\begin{frame}{WindowSHAP --- Nayebi et al. (J. Biomedical Informatics, 2023)}
  \begin{columns}
    \begin{column}{0.51\textwidth}
    	\vspace{-0.3cm}
      \begin{tcolorbox}[
        colback=blue!3, colframe=blue!40!black,
        arc=2mm, boxrule=0.5pt, fonttitle=\bfseries,
        title={Idée centrale}]
        Réduire le coût de KernelSHAP en calculant des valeurs Shapley
        par \textbf{fenêtre temporelle} plutôt que par pas de temps.
      \end{tcolorbox}
      \begin{tcolorbox}[
        colback=orange!3, colframe=orange!50!black,
        arc=2mm, boxrule=0.5pt, fonttitle=\bfseries,
        title={3 variantes}]
        \begin{itemize}
          \item \textbf{Stationary} : fenêtres fixes, non chevauchantes
          \item \textbf{Sliding} : fenêtres chevauchantes (stride $s$)
          \item \textbf{Dynamic} : longueur variable\\
                (divise les fenêtres importantes, agrège les autres)
        \end{itemize}
      \end{tcolorbox}
    \end{column}
    \begin{column}{0.48\textwidth}
    	\vspace{-0.7cm}
      \begin{tcolorbox}[
        colback=green!3, colframe=green!40!black,
        arc=2mm, boxrule=0.5pt, fonttitle=\bfseries,
        title={Résultats clés}]
        \begin{itemize}
          \item \textbf{$-80\%$ CPU} vs KernelSHAP\\
                (10 pas fusionnés, $L=120$)
          \item Dynamic : meilleure qualité car \textbf{pas d'hypothèse} sur les instants importants
          \item 3 datasets cliniques (TBI, MIMIC-III)
        \end{itemize}
      \end{tcolorbox}
      \begin{tcolorbox}[
        colback=red!3, colframe=red!60!black,
        arc=2mm, boxrule=0.5pt, fonttitle=\bfseries,
        title={Limite}]
        Validé uniquement en \textbf{classification}.\\
        Pas de librairie Python officielle.
      \end{tcolorbox}
    \end{column}
  \end{columns}
  \source{Nayebi et al., J. Biomed. Inform. 2023 --- DOI:10.1016/j.jbi.2023.104438}
\end{frame}

%% 4.6a — TsSHAP
\begin{frame}{TsSHAP --- Raykar et al. (IBM Research, 2023) \hfill \textcolor{blue!70!black}{\textbf{[Le plus pertinent]}}}
	\textbf{Problème résolu}
	\rule{\linewidth}{0.5pt}
	\vspace{2pt}
	
	\textbf{Seule méthode SHAP conçue nativement pour la prévision}.\\
	LIME, KernelSHAP, TimeSHAP, WindowSHAP expliquent des scores
	de classification --- pas des prévisions.
	
	\vspace{8pt}
	
	\textbf{Pipeline}
	\rule{\linewidth}{0.5pt}
	\vspace{2pt}
	
	\begin{enumerate}
		\item \textbf{Backtesting} $\to$ prévisions historiques $\hat{y}(T+h|T)$
		\item \textbf{Features interprétables} : lags, saisonnalités, date/heure
		\item \textbf{Surrogate XGBoost} : imite le forecaster
		\item \textbf{TreeSHAP} : valeurs Shapley sur le surrogate
	\end{enumerate}
	
	\vspace{4pt}
	\source{Raykar et al., IBM Research 2023 --- arXiv:2303.12316}
\end{frame}

\begin{frame}{TsSHAP --- Features interprétables}
	\textbf{3 portées d'explication}
	\rule{\linewidth}{0.5pt}
	
	\textbf{Locale} : une prédiction \hfill
	\textbf{Semi-locale} : un horizon \hfill
	\textbf{Globale} : toute la série
	
	\vspace{18pt}
	
	\textbf{Features interprétables (Tables 1-2)}
	\rule{\linewidth}{0.5pt}
	
	\begin{itemize}
		\item \texttt{LagFeatures} : $y(t-k)$
		\item \texttt{SeasonalLagFeatures} : $y(t-k \cdot m)$
		\item \texttt{RollingWindowFeatures} : max, min, mean
		\item \texttt{DateFeatures} : mois, jour, semaine\ldots
		\item \texttt{HolidayFeatures} : jours fériés
	\end{itemize}
	
	\vspace{6pt}
	
	\rule{\linewidth}{0.5pt}
	
	\vspace{6pt}
	
	TsSHAP reste la \textbf{référence} pour l'explicabilité en prévision,
	mais nécessite un surrogate XGBoost (fidélité approximative).
	
	\source{Raykar et al., IBM Research 2023 --- arXiv:2303.12316}
\end{frame}

%% 4.6b — TsSHAP résultats + limites
\begin{frame}{TsSHAP --- Résultats et limites}
  \begin{columns}
    \begin{column}{0.538\textwidth}
      \begin{tcolorbox}[
        colback=blue!3, colframe=blue!40!black,
        arc=2mm, boxrule=0.5pt, fonttitle=\bfseries,
        title={Résultats (Fidélité --- Table 6)}]
        Fidélité moyenne sur 5 datasets :
        \begin{itemize}
          \item \textbf{MovingAverage} : 0,244 (global) / 0,654 (local)
          \item Explications \textbf{semi-locales} toujours plus fidèles
          \item Complexité stable entre portées
        \end{itemize}
      \end{tcolorbox}
      \begin{tcolorbox}[
        colback=orange!3, colframe=orange!50!black,
        arc=2mm, boxrule=0.5pt, fonttitle=\bfseries,
        title={Illustration (jeans-sales-weekly)}]
        \begin{itemize}
          \item SeasonalNaive $\to$ \texttt{sales(t-52)}
          \item Prophet $\to$ \texttt{discount(t)} = principale feature externe
          \item XGBoost $\to$ \texttt{sales(t-1)} domine
        \end{itemize}
      \end{tcolorbox}
    \end{column}
    \begin{column}{0.46\textwidth}
      \begin{tcolorbox}[
        colback=red!3, colframe=red!60!black,
        arc=2mm, boxrule=0.5pt, fonttitle=\bfseries,
        title={Limites}]
        \begin{itemize}
          \item \textbf{Univarié uniquement}\\
                FAYAM = 18 fonctions $\to$ à appliquer par série
          \item \textbf{Surrogate $\neq$ modèle réel}\\
                XGBoost peut mal imiter un Transformer
          \item \textbf{Features à définir manuellement}
        \end{itemize}
      \end{tcolorbox}
      \begin{tcolorbox}[
        colback=green!3, colframe=green!40!black,
        arc=2mm, boxrule=0.5pt, fonttitle=\bfseries,
        title={Pour notre projet}]
        TsSHAP est le \textbf{H2 prioritaire} si H1 (SoftCAM-Transformer) échoue
      \end{tcolorbox}
    \end{column}
  \end{columns}
  \source{Raykar et al., 2023}
\end{frame}

%% 4.7 — SHAPformer
\begin{frame}{SHAPformer --- Hertel et al. (KIT, déc. 2025)}
  \begin{columns}
    \begin{column}{0.518\textwidth}
      \begin{tcolorbox}[
        colback=blue!3, colframe=blue!40!black,
        arc=2mm, boxrule=0.5pt, fonttitle=\bfseries,
        title={Problème résolu}]
        Calculer des valeurs SHAP \textbf{exactes} pour les Transformers de prévision
        \textbf{sans échantillonnage} --- en moins d'une seconde.
      \end{tcolorbox}
      \begin{tcolorbox}[
        colback=orange!3, colframe=orange!50!black,
        arc=2mm, boxrule=0.5pt, fonttitle=\bfseries,
        title={Mécanisme}]
        \begin{enumerate}
          \item Grouper les features en $N$ groupes
          \item Entraîner avec \textbf{entrées masquées} (masked attention)
          \item Inférence : évaluer sur les $2^N$ sous-ensembles via attention
        \end{enumerate}
        $\Rightarrow$ SHAP \textbf{exact}, pas de Monte Carlo
      \end{tcolorbox}
    \end{column}
    \begin{column}{0.48\textwidth}
      \begin{tcolorbox}[
        colback=green!3, colframe=green!40!black,
        arc=2mm, boxrule=0.5pt, fonttitle=\bfseries,
        title={Résultats (TransnetBW, Table 1)}]
        \hspace{-0.5cm}
        \begin{tabular}{lc}
          \toprule
          \textbf{Méthode} & \textbf{Inférence} \\
          \midrule
          Permutation Explainer & 484 s \\
          Custom Masker & 3,54 s \\
          \textbf{SHAPformer} & \textbf{0,60 s} \\
          \bottomrule
        \end{tabular}
        \vspace{4pt}
        \textbf{800$\times$ plus rapide} | RMSE comparable
      \end{tcolorbox}
      \begin{tcolorbox}[
        colback=red!3, colframe=red!60!black,
        arc=2mm, boxrule=0.5pt, fonttitle=\bfseries,
        title={Limite}]
        Réentraînement requis.\\
        Complexité $\propto 2^N$ groupes de features.
      \end{tcolorbox}
    \end{column}
  \end{columns}
  \source{Hertel et al., arXiv:2512.20514, déc. 2025 --- \href{https://github.com/KIT-IAI/SHAPformer}{github.com/KIT-IAI/SHAPformer}}
\end{frame}


%% ------------------------------------------------------------------
%%  SECTION 5 — FAMILLE CAM
%% ------------------------------------------------------------------
\section{Famille CAM : de l'activation à l'évidence}

%% 5.1 — CAM
\begin{frame}{CAM --- Zhou et al. (CVPR 2016) : l'origine}
  \begin{columns}
    \begin{column}{0.55\textwidth}
      \begin{tcolorbox}[
        colback=blue!3, colframe=blue!40!black,
        arc=2mm, boxrule=0.5pt, fonttitle=\bfseries,
        title={Class Activation Maps}]
        Pour un CNN classifieur :
        \[
          M_c(x,y) = \sum_k w_k^c \cdot f_k(x,y)
        \]
        \vspace{-0.8cm}
        \begin{itemize}
          \item $f_k$ : $k$-ème feature map dernière couche conv.
          \item $w_k^c$ : poids FC pour la classe $c$
          \item $M_c$ : carte de chaleur des régions discriminantes
        \end{itemize}
      \end{tcolorbox}
      \begin{tcolorbox}[
        colback=red!3, colframe=red!60!black,
        arc=2mm, boxrule=0.5pt, fonttitle=\bfseries,
        title={Contrainte}]
        Requiert un GAP avant la couche FC.\\
        Architecture imposée, pas universelle.
      \end{tcolorbox}
    \end{column}
    \begin{column}{0.41\textwidth}
      \begin{tcolorbox}[
        colback=orange!3, colframe=orange!50!black,
        arc=2mm, boxrule=0.5pt, fonttitle=\bfseries,
        title={Principe visuel}]
        \begin{itemize}
          \item Image $\to$ CNN $\to$ Feature maps
          \item GAP $\to$ FC $\to$ Classe
          \item Carte : quelles régions activent la décision ?
        \end{itemize}
      \end{tcolorbox}
      \begin{tcolorbox}[
        colback=green!3, colframe=green!40!black,
        arc=2mm, boxrule=0.5pt, fonttitle=\bfseries,
        title={Intérêt}]
        Première méthode à \textbf{localiser} ce que le modèle regarde,
        sans back-propagation ni perturbation
      \end{tcolorbox}
    \end{column}
  \end{columns}
  \source{Zhou et al., CVPR 2016}
\end{frame}

%% 5.2 — GradCAM
\begin{frame}{GradCAM --- Selvaraju et al. (ICCV 2017) : extension gradient}
  \begin{columns}
    \begin{column}{0.435\textwidth}
      \begin{tcolorbox}[
        colback=blue!3, colframe=blue!40!black,
        arc=2mm, boxrule=0.5pt, fonttitle=\bfseries,
        title={Contribution}]
        Généralise CAM à toute architecture CNN sans contrainte GAP,
        via les \textbf{gradients} du score de classe $c$ :
        \[
          \alpha_k^c = \frac{1}{Z} \sum_{i,j} \frac{\partial y^c}{\partial A^k_{ij}}
        \]
        \[
          L^c_{\text{GradCAM}} = \text{ReLU}\!\left(\sum_k \alpha_k^c A^k\right)
        \]
      \end{tcolorbox}
    \end{column}
    \begin{column}{0.545\textwidth}
      \begin{tcolorbox}[
        colback=green!3, colframe=green!40!black,
        arc=2mm, boxrule=0.5pt, fonttitle=\bfseries,
        title={Avantages vs CAM}]
        \begin{itemize}
          \item Applicable à \textbf{toute couche} convolutive
          \item Pas de contrainte architecturale
          \item Cartes de résolution variable
        \end{itemize}
      \end{tcolorbox}
      \begin{tcolorbox}[
        colback=red!3, colframe=red!60!black,
        arc=2mm, boxrule=0.5pt, fonttitle=\bfseries,
        title={Limites persistantes}]
        \begin{itemize}
          \item Méthode \textbf{post-hoc} : un back-pass par classe
          \item Résolution limitée à la feature map cible
          \item Sensible aux gradients saturés
        \end{itemize}
      \end{tcolorbox}
    \end{column}
  \end{columns}
  \source{Selvaraju et al., ICCV 2017}
\end{frame}

%% 5.3a — SoftCAM architecture
\begin{frame}{SoftCAM --- Djoumessi \& Berens (arXiv, mai 2025)}
	\textbf{Problème fondamental posé par les auteurs}
	\rule{\linewidth}{0.5pt}
	
	Les méthodes post-hoc (GradCAM, SHAP, IG\ldots) approximent le raisonnement \textit{après coup}.
	Elles sont \textbf{sensibles aux perturbations}, \textbf{infidèles} et ne reflètent pas
	le vrai processus décisionnel.
	\vspace{-0.5pt}
	\rule{\linewidth}{0.5pt}
	
	\vspace{4pt}
	
	\begin{minipage}{0.46\textwidth}
		\textbf{Architecture SoftCAM}
		\rule{\linewidth}{0.5pt}
		
		{\footnotesize
			\textbf{CNN standard :} \\
			$\text{Backbone} \to \text{GAP} \to \text{FCL} \to \text{Softmax}$
			
			\vspace{4pt}
			\textbf{SoftCAM (modification minimale) :} \\
			$\text{Backbone} \to \underbrace{\text{Conv}_{1\times1}}_{\text{class-evidence layer}} \to \text{AvgPool} \to \text{Softmax}$
			
			\textbf{Sortie :} $A \in \mathbb{R}^{M \times N \times C}$ --- \textbf{carte d'évidence}
		}
	\end{minipage}
	\hfill
	\begin{minipage}{0.51\textwidth}
		\textbf{Propriété clé}
		\rule{\linewidth}{0.5pt}
		{\small
			\begin{itemize}
				\item \textbf{1 seul forward pass} pour toutes les classes
				\item Explication \textbf{intrinsèque} : la carte \textit{est} la prédiction
				\item Aucun paramètre supplémentaire vs FCL
			\end{itemize}
		}
	\end{minipage}
	
	\vspace{-0.3pt}
	\rule{\linewidth}{0.5pt}
	
	
	\textbf{Avant / Après}
	\vspace{-0.5pt}
	\rule{\linewidth}{0.5pt}
	
	{\small
		Post-hoc $\to$ back-pass ou perturbations \\
		SoftCAM $\to$ \textbf{0 coût supplémentaire}
	}
	\source{Djoumessi \& Berens, arXiv:2505.17748, mai 2025}
\end{frame}

%% 5.3b — SoftCAM régularisation
\begin{frame}{SoftCAM --- Régularisation ElasticNet}
  \begin{tcolorbox}[
    colback=blue!3, colframe=blue!40!black,
    arc=2mm, boxrule=0.5pt, fonttitle=\bfseries,
    title={Loss augmentée (Eq. 5)}]
    \[
      \mathcal{L}(y, \hat{y}) = \underbrace{\text{CE}(y, \hat{y})}_{\text{performance}}
      + \underbrace{\lambda_1 \sum_{c,i,j} |A_c^{i,j}|}_{\ell_1 \text{ --- Lasso}}
      + \underbrace{\lambda_2 \sum_{c,i,j} \|A_c^{i,j}\|_2^2}_{\ell_2 \text{ --- Ridge}}
    \]
  \end{tcolorbox}

  \begin{columns}
    \begin{column}{0.31\textwidth}
      \begin{tcolorbox}[
        colback=orange!3, colframe=orange!50!black,
        arc=2mm, boxrule=0.5pt, fonttitle=\bfseries,
        title={Sparse SoftCAM ($\lambda_1 > 0$)}, halign title=center]
        \centering
        Lasso $\to$ cartes \textbf{parcimonieuses}\\[4pt]
        Lésions focales, régions précises
      \end{tcolorbox}
    \end{column}
    \begin{column}{0.34\textwidth}
      \begin{tcolorbox}[
        colback=blue!3, colframe=blue!40!black,
        arc=2mm, boxrule=0.5pt, fonttitle=\bfseries,
        title={Dense SoftCAM ($\lambda_2 > 0$)}, halign title=center]
        \centering
        Ridge $\to$ cartes \textbf{lissées}\\[4pt]
        Grandes régions, faux négatifs coûteux
      \end{tcolorbox}
    \end{column}
    \begin{column}{0.31\textwidth}
      \begin{tcolorbox}[
        colback=green!3, colframe=green!40!black,
        arc=2mm, boxrule=0.5pt, fonttitle=\bfseries,
        title={ElasticNet ($\lambda_1 + \lambda_2$)}, halign title=center]
        \centering
        \textbf{Compromis} adaptatif\\[4pt]
        Trade-off contrôlable par $\lambda$
      \end{tcolorbox}
    \end{column}
  \end{columns}

  \source{Djoumessi \& Berens, 2025 --- $\lambda_1 \in [10^{-6}, 9\times10^{-4}]$, $\lambda_2 \in [7\times10^{-5}, 2\times10^{-4}]$}
\end{frame}

%% 5.3c — SoftCAM résultats
\begin{frame}{SoftCAM --- Résultats expérimentaux}
	\begin{minipage}{0.52\textwidth}
		\textbf{Performance préservée (Table 1)}
		\rule{\linewidth}{0.5pt}
		\vspace{2pt}
		
		{\small
			\begin{tabular}{lcc}
				\toprule
				\textbf{Modèle} & \textbf{Backbone} & \textbf{AUC} \\
				\midrule
				Standard & VGG-16 & 0,938 \\
				SoftCAM dense & VGG-16 & \textbf{0,942} \\
				SoftCAM sparse & VGG-16 & 0,938 \\
				Standard & ResNet-50 & 0,999 \\
				SoftCAM & ResNet-50 & \textbf{1,000} \\
				\bottomrule
		\end{tabular}}
		
		\vspace{4pt}
		{\small \textit{La modification n'altère pas la précision}}
	\end{minipage}
	\hfill
	\begin{minipage}{0.44\textwidth}
		\vspace{-1.05cm}
		\textbf{Explicabilité (Fig. 3)}
		\rule{\linewidth}{0.5pt}
		\vspace{2pt}
		
		{\footnotesize
			Sparse SoftCAM ResNet \textbf{surpasse toutes les méthodes post-hoc}
			en top-k precision (fundus + OCT) : \\
			vs GradCAM, IG, Guided Backprop, ScoreCAM, LayerCAM}
	\end{minipage}
	
	\vspace{14pt}
	
	\textbf{Coût computationnel}
	\rule{\linewidth}{0.5pt}
	
	{\footnotesize
		\textbf{1 forward pass} contre : 1 back-pass (gradient-based) ou $N$ forward passes (perturbation-based)}
	
	\vspace{4pt}
	\source{Djoumessi \& Berens, 2025 --- 3 datasets médicaux : Kaggle DR Fundus, Retinal OCT, RSNA CXR}
\end{frame}

%% 5.3d — SoftCAM perspective ViT + lien H1
\begin{frame}{SoftCAM --- Perspective et lien avec notre projet}
  \begin{columns}
    \begin{column}{0.50\textwidth}
      \begin{tcolorbox}[
        colback=red!3, colframe=red!60!black,
        arc=2mm, boxrule=0.5pt, fonttitle=\bfseries,
        title={Limites identifiées}]
        \begin{itemize}
          \item Résolution coarse (feature maps 16$\times$16)
          \item Évalué uniquement sur CNN (imagerie médicale 2D)
          \item Trade-off $\lambda$ task-specific (à calibrer)
        \end{itemize}
      \end{tcolorbox}
      \begin{tcolorbox}[
        colback=blue!3, colframe=blue!40!black,
        arc=2mm, boxrule=0.5pt, fonttitle=\bfseries,
        title={Porte ouverte (p. 10)}]
        \textit{``In the future, we could explore the integration of SoftCAM
        with other standard architectures like ViT.''}\\[4pt]
        $\Rightarrow$ \textbf{C'est exactement notre H1}
      \end{tcolorbox}
    \end{column}
    \begin{column}{0.46\textwidth}
      \begin{tcolorbox}[
        colback=green!3, colframe=green!40!black,
        arc=2mm, boxrule=0.5pt, fonttitle=\bfseries,
        title={Notre transposition --- H1}]
        \textbf{Adapter SoftCAM au TimeSeriesTransformer de FAYAM} :
        \begin{itemize}
          \item Couche cible : \textbf{projection finale du décodeur}\\
                ($d_{\text{model}} \to$ distribution)
          \item Evidence map : $[\text{predictionLength} \times \text{contextLength}]$
          \item ElasticNet : sparsité temporelle (lags critiques)
          \item Validation : comparer avec attention maps (H3)
        \end{itemize}
      \end{tcolorbox}
    \end{column}
  \end{columns}
\end{frame}


%% ------------------------------------------------------------------
%%  SECTION 6 — COMPARAISON SYNTHÉTIQUE
%% ------------------------------------------------------------------
\section{Comparaison synthétique}

%% 6.1
\begin{frame}{Tableau comparatif}
	\begin{tcolorbox}[colback=blue!3, colframe=blue!40!black, arc=2mm, boxrule=0.5pt]
		\tiny
		\begin{tabular}{lccccccc}
			\toprule
			\textbf{Méthode} & \textbf{Famille} & \textbf{Paradigme} & \textbf{Agnostic} & \textbf{Prévision} & \textbf{Transformer} & \textbf{Fidélité} & \textbf{Coût inférence} \\
			\midrule
			LIME        & Perturbation & Post-hoc & \checkmark & \texttimes & \checkmark & Faible  & $N$ evals \\
			KernelSHAP  & SHAP & Post-hoc & \checkmark & \texttimes & \checkmark & Bonne   & $\sim 2^{|F|}$ evals \\
			TimeSHAP    & SHAP & Post-hoc & \checkmark & \texttimes & partiel    & Bonne   & $O(l \cdot 2^{l-i})$ \\
			WindowSHAP  & SHAP & Post-hoc & \checkmark & \texttimes & partiel    & Bonne   & $-80\%$ CPU \\
			TsSHAP      & SHAP & Post-hoc & \checkmark & \checkmark & \checkmark & Bonne   & TreeSHAP rapide \\
			SHAPformer  & SHAP & Post-hoc* & partiel   & \checkmark & \checkmark & Exacte  & $<1$ s \\
			\midrule
			CAM         & CAM & Post-hoc & \texttimes & \texttimes & \texttimes & Mod.    & 1 back-pass \\
			GradCAM     & CAM & Post-hoc & \texttimes & \texttimes & \texttimes & Mod.    & 1 back-pass \\
			\textbf{SoftCAM} & \textbf{CAM} & \textbf{Intrinsèque} & \texttimes & \textbf{\checkmark*} & \textbf{\checkmark*} & \textbf{Maximale} & \textbf{1 fwd pass} \\
			\bottomrule
		\end{tabular}
	\end{tcolorbox}
	
	\vspace{2pt}
	\tiny
	\textcolor{gray}{*} Extension proposée dans notre projet (H1) \hfill * SHAPformer requiert réentraînement (pas modification architecturale)
\end{frame}

%% 6.2
\begin{frame}{Lecture du tableau --- Aucune méthode n'est parfaite pour notre cas}
  \begin{columns}
    \begin{column}{0.53\textwidth}
      \begin{tcolorbox}[
        colback=red!3, colframe=red!60!black,
        arc=2mm, boxrule=0.5pt, fonttitle=\bfseries,
        title={Ce que le tableau révèle}]
        \begin{itemize}
          \item LIME / KernelSHAP : pas de séries temporelles
          \item TimeSHAP / WindowSHAP : pas de prévision
          \item TsSHAP : prévision \checkmark, mais univarié + surrogate
          \item SHAPformer : Transformer + prévision \checkmark, mais réentraînement
          \item SoftCAM original : CNN uniquement
        \end{itemize}
      \end{tcolorbox}
    \end{column}
    \begin{column}{0.43\textwidth}
      \begin{tcolorbox}[
        colback=green!3, colframe=green!40!black,
        arc=2mm, boxrule=0.5pt, fonttitle=\bfseries,
        title={Ce qui manque dans la littérature}]
        Une méthode \textbf{simultanément} :
        \begin{itemize}
          \item[\checkmark] Adaptée aux \textbf{Transformers}
          \item[\checkmark] Conçue pour la \textbf{prévision}
          \item[\checkmark] À \textbf{fidélité intrinsèque}
          \item[\checkmark] Séries \textbf{multivariées}
          \item[\checkmark] Sans surcoût à l'inférence
        \end{itemize}
      \end{tcolorbox}
    \end{column}
  \end{columns}

  \vspace{4pt}
    \centering
    \textbf{C'est précisément ce vide que notre H1 cherche à combler}

\end{frame}


%% ------------------------------------------------------------------
%%  SECTION 7 — RETOUR SUR FAYAM : QUELLE VOIE CHOISIR ?
%% ------------------------------------------------------------------
\section{Retour sur FAYAM : quelle voie choisir ?}

%% 7.1
\begin{frame}{Architecture TimeSeriesTransformer --- La couche cible}
  \begin{tcolorbox}[
    colback=blue!3, colframe=blue!40!black,
    arc=2mm, boxrule=0.5pt, fonttitle=\bfseries,
    title={Architecture \texttt{TimeSeriesTransformer} HuggingFace}]
    \begin{center}
    \begin{tikzpicture}[scale=0.70, transform shape,
      box/.style={draw, rounded corners, fill=fayam!15, text width=3cm, align=center, minimum height=0.8cm}]
      \node[box] (enc) {Encodeur\\(contexte passé)};
      \node[box, right=1cm of enc] (dec) {Décodeur\\(horizon futur)};
      \node[box, right=1cm of dec, fill=verrou!20] (proj) {\textbf{Projection finale}\\$d_\text{model} \to$ distribution\\{\small \textit{couche cible H1}}};
      \node[box, right=1cm of proj, fill=softcam!20] (out) {Prédiction\\$\hat{y}(T+h)$};
      \draw[-{Stealth}] (enc) -- (dec);
      \draw[-{Stealth}] (dec) -- (proj);
      \draw[-{Stealth}] (proj) -- (out);
    \end{tikzpicture}
    \end{center}
  \end{tcolorbox}

  \begin{columns}
    \begin{column}{0.50\textwidth}
      \begin{tcolorbox}[
        colback=green!3, colframe=green!40!black,
        arc=2mm, boxrule=0.5pt, fonttitle=\bfseries,
        title={Adaptation SoftCAM (H1)}]
        Remplacer la projection finale par une \textbf{couche d'évidence temporelle}
        produisant un tenseur $[\text{predictionLength} \times \text{contextLength}]$
      \end{tcolorbox}
    \end{column}
    \begin{column}{0.46\textwidth}
      \begin{tcolorbox}[
        colback=red!3, colframe=red!60!black,
        arc=2mm, boxrule=0.5pt, fonttitle=\bfseries,
        title={Risque H1}]
        La convolution 1$\times$1 préserve les positions spatiales d'un CNN.
        Un Transformer mélange les positions par self-attention.
        L'extension n'est pas triviale.
      \end{tcolorbox}
    \end{column}
  \end{columns}
\end{frame}

%% 7.2
\begin{frame}{Notre grille de décision}
	\textbf{H1 --- SoftCAM-Transformer (prioritaire)}
	\rule{\linewidth}{0.5pt}
	
	\begin{itemize}
		\item Modifier la projection finale du décodeur
		\item Loss = CE + ElasticNet sur la carte d'évidence
		\item Réentraîner le Transformer avec cette contrainte
		\item Résultat : carte $[\text{predL} \times \text{ctxL}]$ intrinsèque
	\end{itemize}
	
	\vspace{2pt}
	\textbf{Avantage} : fidélité maximale, 1 forward pass\\
	\textbf{Risque} : incompatibilité CNN $\to$ Transformer
	
	\vspace{8pt}
	
	\textbf{H2 --- TsSHAP / SHAPformer (repli)}
	\rule{\linewidth}{0.5pt}
	
	\begin{itemize}
		\item TsSHAP : surrogate XGBoost sur backtests FAYAM
		\item SHAPformer : réentraîner avec masked attention
	\end{itemize}
	
	\vspace{2pt}
	\textbf{Avantage} : sans modifier l'architecture (TsSHAP) ou code public (SHAPformer)\\
	\textbf{Risque} : fidélité moindre (surrogate)
\end{frame}

%% 7.3
\begin{frame}{Notre hypothèse prioritaire : H1 --- SoftCAM-Transformer}
  \begin{tcolorbox}[
    colback=blue!3, colframe=blue!40!black,
    arc=2mm, boxrule=0.5pt, fonttitle=\bfseries,
    title={Contribution scientifique visée}]
    \centering
    \textbf{Première transposition documentée d'un schéma SoftCAM-like\\
    à un Transformer de séries temporelles},\\
    avec étude différentielle par cluster DBSCAN (profils de charge FaaS)
  \end{tcolorbox}

  \vspace{6pt}
  \begin{columns}
    \begin{column}{0.52\textwidth}
      \begin{tcolorbox}[
        colback=orange!3, colframe=orange!50!black,
        arc=2mm, boxrule=0.5pt, fonttitle=\bfseries,
        title={Critère de bascule vers H2}]
        Si après un bon moment :
        \begin{itemize}
          \item L'adaptation SoftCAM ne converge pas, \textbf{ou}
          \item La performance de prédiction dégrade significativement
        \end{itemize}
        $\Rightarrow$ \textbf{Basculer sur H2 sans attendre}
      \end{tcolorbox}
    \end{column}
    \begin{column}{0.45\textwidth}
      \begin{tcolorbox}[
        colback=gray!5, colframe=gray!50!black,
        arc=2mm, boxrule=0.5pt, fonttitle=\bfseries,
        title={H3 --- Attention weights (validation)}]
        Si H1 ou H2 aboutit, les attention maps extraites par
        \texttt{output\_attentions=True} servent à \textbf{valider}
        la cohérence des explications obtenues
      \end{tcolorbox}
    \end{column}
  \end{columns}
\end{frame}

%% 7.4
\begin{frame}{Questions ouvertes et prochaines étapes}
	\textbf{Questions méthodologiques ouvertes}
	\rule{\linewidth}{0.5pt}
	
	\begin{enumerate}
		\item Comment définir l'équivalent d'une \textit{class evidence map}
		pour la \textbf{régression} multi-horizons ?
		\item Quel est l'équivalent temporel des super-pixels CNN ?
		(patch temporel ? lag ? cluster DBSCAN ?)
		\item Comment évaluer la fidélité sur une tâche de prévision ?
		(comprehensiveness / sufficiency adaptés)
	\end{enumerate}
	
	
	\textbf{Prochaines étapes (Phase 1)}
	\rule{\linewidth}{0.5pt}
	
	\begin{enumerate}
		\item Récupérer le code FAYAM
		\item Agréger le dataset + EDA
		\item Entraîner \texttt{TimeSeriesTransformer} (\texttt{output\_attentions=True})
		\item Cartographier la projection du décodeur
		\item Métriques FAYAM (sMAPE, RMSE, R², Spearman)
	\end{enumerate}
\end{frame}

\begin{frame}{Questions ?}
	\vspace{12pt}

	\begin{center}
		\Large\textbf{Merci pour votre attention}

		\vspace{12pt}

		\large Des questions ?
	\end{center}
\end{frame}


%% ------------------------------------------------------------------
%%  RÉFÉRENCES
%% ------------------------------------------------------------------
\begin{frame}[allowframebreaks]{Références}
  \scriptsize
  \begin{thebibliography}{99}

    \bibitem{Ribeiro2016}
      M.T. Ribeiro, S. Singh, C. Guestrin.
      \textit{``Why Should I Trust You?'' Explaining the Predictions of Any Classifier}.
      KDD 2016. arXiv:1602.04938.

    \bibitem{Lundberg2017}
      S.M. Lundberg, S.-I. Lee.
      \textit{A Unified Approach to Interpreting Model Predictions}.
      NeurIPS 2017.

    \bibitem{Bento2021}
      J. Bento, P. Saleiro, A.F. Cruz, M.A.T. Figueiredo, P. Bizarro.
      \textit{TimeSHAP: Explaining Recurrent Models through Sequence Perturbations}.
      KDD 2021. arXiv:2012.00073.

    \bibitem{Nayebi2023}
      A. Nayebi, M. Tipirneni, C. Reddy, F. Foreman, V. Subbian.
      \textit{WindowSHAP: An Efficient Framework for Explaining Time-Series Classifiers
      Based on Shapley Values}.
      Journal of Biomedical Informatics, 2023. DOI:10.1016/j.jbi.2023.104438.

    \bibitem{Raykar2023}
      V.C. Raykar, A. Jati, S. Mukherjee, N. Aggarwal, K. Sarpatwar,
      G. Ganapavarapu, R. Vaculin.
      \textit{TsSHAP: Robust Model Agnostic Feature-Based Explainability for
      Time Series Forecasting}.
      arXiv:2303.12316, IBM Research, 2023.

    \bibitem{Hertel2025}
      M. Hertel, S. Pütz, R. Mikut, V. Hagenmeyer, B. Schäfer.
      \textit{SHAPformer: Explainable Time Series Forecasting with Sampling-Free
      SHAP Values for Transformers}.
      arXiv:2512.20514, KIT, décembre 2025.

    \bibitem{Zhou2016}
      B. Zhou, A. Khosla, A. Lapedriza, A. Oliva, A. Torralba.
      \textit{Learning Deep Features for Discriminative Localization}.
      CVPR 2016.

    \bibitem{Selvaraju2017}
      R.R. Selvaraju, M. Cogswell, A. Das, R. Vedantam, D. Parikh, D. Batra.
      \textit{Grad-CAM: Visual Explanations from Deep Networks via
      Gradient-Based Localization}.
      ICCV 2017.

    \bibitem{Djoumessi2025}
      K. Djoumessi, P. Berens.
      \textit{SoftCAM: Making Black Box Models Self-Explainable}.
      arXiv:2505.17748, mai 2025.

    \bibitem{FAYAM2024}
      J. Fayam \textit{et al.}
      \textit{Prédiction de charge FaaS par réseaux de neurones profonds}.
      Mémoire Master 2, Université de Dschang, 2024.

  \end{thebibliography}
\end{frame}

%% ============================================================
\end{document}
%% ============================================================
